\section{Proof obligations}
\label{sec:obligations}

This section keeps two ledgers. The first, \cref{sec:obligations-current},
holds the obligations open against the current specification, CCME
3.1.0-rc.2 at repository commit \code{252216ba}, and against the reference
implementation at the same commit: places where the normative text
contradicts itself, where the text states a boundary of a guarantee rather
than a rule, and where the engine departs from the text. Each states the
claim, a minimal scenario, the consequence for the results of this report, the
repair or the recorded status, and the case and test that pin it. They are
ordered by severity, the silent loss of a release first and diagnostic
metadata last. The second ledger, \cref{sec:obligations-historical}, is the
audit of CCME 1.0.0 at \code{713f1a0b} that this report was first written for.
Its six obligations are provenance: the released 2.0.0 text incorporated
their repairs and the current text keeps them, so they are reproduced
unchanged.

Every \SPEC{} reference in \cref{sec:obligations-current} is to the current
text. In \cref{sec:obligations-historical}, except where a paragraph
explicitly names CCME 2.0.0 or a current theorem, every numbered reference
and quoted guarantee refers to the 1.0.0 snapshot, and present tense
describes the historical rule being executed.

\subsection{Current obligations}
\label{sec:obligations-current}

The sources consulted are \SPEC{} itself, its dated design history and
changelog, the release notes of the engine, the integration test plan, and
the test sources of the engine and of its integration suite. Where the design
history names a departure, the obligation quotes it; where it does not, the
obligation says what the code does and where.

\subsubsection{The owed window and \diag{E201} are specified but not implemented}

\begin{obligation}[a consumer that sits out the run in which its provider delivers is stranded]
\label{ob:owed}
\spec{13.4a} admits a bump-axis contribution for a target until the source
has \emph{delivered} it. \spec{13.3} adds, for every edge $D \to P$ where $D$
has a baseline, an owed window \code{owedWindow(P, D)} to the union of
windows parsed in \spec{13.4}, so that a debt survives the release of the
provider; and \spec{19.3} refuses with \diag{E201} a run that would release a
provider at the baseline commit of a consumer it still owes unless the
consumer releases after it in the same run. The reference implementation
implements the delivery test and the reconciliation of a proceeding consumer,
and neither the owed window nor \diag{E201}.
\end{obligation}

\paragraph{Source.} The design history entry of 22 September 2026,
``Delivery discharges a propagated contribution'', ends: ``The dispat engine
implements the delivery admission and the reconciliation of a proceeding
consumer; the owed windows and \code{E201} are not yet implemented, and a
consumer that sits out the run in which its provider releases the owed commit
is therefore still stranded in that engine, which its release notes list as a
departure.'' The engine's sources at the same commit contain no \code{E201}
and no owed window; \code{admission.go} carries \code{isDelivered} and
\code{owedSources}, the delivery half. The release notes at this commit do
not yet carry the departure: \code{services/dispat/CHANGELOG.md} ends at
\code{v1.11.0-rc.3} of 20 September, two days before the entry, and
\code{announce/rc/notes.yaml} does not mention it.

\paragraph{Scenario.} One commit $C_1$ carries \code{feat(core)\^{}: x} and
\code{feat(cli): y}, with \code{cli} consuming \code{core}. Run 1:
\code{core} fails to publish; \code{cli} has a cause of its own, proceeds at
\code{2.1.0} with its manifest naming \code{core}'s baseline \code{1.4.0},
and is tagged at $C_1$ (\spec{19.3}; vector 80d). A commit $C_2$ lands. Run 2
releases \code{core@1.5.0} at $C_2$ in a run \code{cli} sits out, held or
deselected. Now $W(\code{core}) = \emptyset$ and $W(\code{cli}) = \{C_2\}$:
$C_1$ lies in no window of either package, so the unit is never parsed,
\code{cli} is never planned, and its manifest names \code{core@1.4.0} for
ever. The delivery test would answer correctly if it were asked, since
\code{core}'s release commit $C_2$ is not reached by \code{cli}'s baseline
$C_1$; the owed window is what asks it. Vector 80d states the case in its
last sentence: ``\code{core} published at a later commit in a run \code{cli}
sat out leaves \code{C1} in \code{cli}'s owed window (\S13.3), so the next run
plans \code{cli}'s catch-up.''

\paragraph{Consequence.} In the specification, \cref{thm:g2} holds as
stated: the owed window keeps the contribution admissible until \code{cli}
releases at or after \code{core}'s release carrying $C_1$. In the engine, the
theorem holds only while some pending window still holds the commit, which
is every case but this one. The loss is silent in the plan. The tag-level
screen of \cref{prop:screen} does fire, because \code{core}'s stable commit
is not an ancestor of \code{cli}'s baseline, so a scheduled \spec{13.7b}
audit reports \diag{W195}; nothing releases it.

\paragraph{Repair.} Add the owed windows to the parsed union and \diag{E201}
to the publish protocol, as written. Until then the operator-side remedy is
the one \spec{19.3} states for the refused run: release the consumer in the
same run as the provider, or pin it with an explicit \code{Release-As}.

\paragraph{Pin.} Vector 80d, last two cases. The tests
\code{TestAdmissionCatchesUpAConsumerThatProceededOnItsOwn} and
\code{TestAdmissionCatchesUpAConsumerThatOvertookAHeldProvider} of the
integration suite and
\code{TestAdmissionCatchesUpAConsumerThatOvertookItsProvider} of
\code{internal/plan} pin the delivery half at an unchanged \Head{}; no test
in either suite exercises the run the consumer sits out.

\subsubsection{Reached-set attribution in the measured build}

\begin{obligation}[a unit's whole source set was credited to every reached dependent]
\label{ob:attribution}
\spec{9.2} attributes to a dependent $d$ the sources of a unit within the
unit's depth of $d$, and \spec{13.4a} draws the owed set from that reaching
set. The build that the campaigns of \cref{sec:evaluation} ran,
\code{1.11.0-rc.4+main.272772cb}, credited a unit's whole source set to every
dependent the walk reached: for a unit written over several packages, a
consumer two edges from one source and one from another was owed by both,
so it was released again on the far source's account after the near source
had delivered; and a consumer of only one of the packages was told it
releases because of the other. This is the (R) shape of
\cref{sec:guarantees} and the second shape of \cref{rem:owed-reach}.
\end{obligation}

\paragraph{Source.} Reported by the product session that ran the campaigns,
on 22 September 2026, as a correction of its own earlier statement that the
engine's attribution matched the specification; this report had not repeated
that statement. The fix is commit \code{e461d7d9} of the product repository's
\code{main}, ``attribute a unit's sources to a dependent within the unit's
depth of it'' (\code{internal/plan/admission.go}, \code{plan.go},
\code{propagate.go} and their tests; \code{owedSources} now answers a subset
of the reaching set in name order), after the snapshot \code{252216ba} this
report transcribes and unpublished when this was written. The commit message
records that \code{TestMultiScopeSourceOutsideTheManifestsStaysOutOfUpdates}
had pinned the old attribution and now expects the specified provenance.

\paragraph{Consequence for the evidence.} None for the figures of
\cref{sec:eval:rq5}: every trigger commit in the kernel and LLVM fixtures
scopes one package, and for a single-source unit the reaching set equals the
whole source set, so plans, tags and timings are identical under both rules,
and the reproduction kit's pin at \code{272772cb} stands. The companion
campaigns through 1.10.0 predate the delivery rule and are not affected.

\paragraph{Status.} Fixed after the measured build; the specification's text
still permits the shape, which \cref{rem:owed-reach} records as a
disagreement between its two formulations rather than an engine matter.

\subsubsection{Two releases on one commit: a stated boundary of completeness}

\begin{obligation}[the same-commit failure case of \diag{E201}]
\label{ob:samecommit}
\spec{19.3} states: ``If the consumer is published after the provider and
fails, the same state arises, and \code{E201} reports it with its remedy: an
explicit \code{Release-As: <version>} on the consumer (\S8.6), because no
later plan can compute the debt.'' The protocol reports this state; it does
not resolve it. Completeness (\cref{thm:g2}) therefore has a boundary the
specification itself states.
\end{obligation}

\paragraph{Scenario.} The state of \cref{ob:owed} after run 1, and run 2 at
the same \Head{} $= C_1$ with both packages in the plan: \code{core@1.5.0}
publishes and is tagged at $C_1$, then \code{cli@2.1.1} fails. Both
\code{core}'s release commit and \code{cli}'s baseline are $C_1$. By
\spec{13.4a}, \code{delivered(core, C1, cli)} asks whether some release
commit $t$ of \code{core} carrying $C_1$ lies in the reach of \code{cli}'s
baseline commit; $t = C_1$ does, so nothing is owed, although \code{cli}
still names \code{core@1.4.0}. \spec{13.4a} says why no rule over tags can do
better: ``Two releases on one commit have no order (\S13.7b), so
\code{delivered} cannot tell a target that released at commit \code{t} after
its source from one that released there in an earlier run whose source then
released at \code{t} too.''

\paragraph{Consequence.} \Cref{thm:g2} discharges a contribution when $d$
releases at a commit whose ancestry contains $\cof{u}$. Ancestry is a partial
order on commits and says nothing about two records on one commit, so the
theorem is exact only under the hypothesis that a provider is never released
at the baseline commit of a consumer it owes unless that consumer releases
after it in the same run and succeeds. \spec{19.3} enforces the first half
before publication and reports the second with a remedy. The closing remark
of \cref{prop:screen}, that this case is prevented rather than detected,
holds with the qualification that the prevention is a refusal and the
residue is a report.

\paragraph{Status.} A boundary, stated. In the engine it is unreported,
because \diag{E201} is not implemented (\cref{ob:owed}); the state arises
exactly as described and leaves no diagnostic. Vector 80d names the refused
run; no vector or test exercises the failure after the provider's
publication.

\subsubsection{The version summary still contradicts the graduation rule}

\begin{obligation}[\spec{13.9} does not preserve \spec{11.5}'s train target]
\label{ob:graduation}
\spec{11.5} states the graduation version as ``\code{applyBump(S, E)}, the
same \code{target} as \S11.4, with no prerelease suffix'', where \spec{11.4}
defines $E$ as ``the effective bump accumulated over the pending window
(\S13.3)'', the train window $W(P)$. \spec{13.9}'s stable branch instead
reads ``\code{next = applyBump(stableBaseline(P), effective(P))}'', and
\code{effective(P)} is \code{max(direct(P), propagated(P))} (\spec{13.7})
with \code{direct(P)} taken ``over surviving tuples for \code{P} whose commit
is in \code{Wfresh(P)}'' (\spec{13.6}). On a channel-only graduation the
fresh effective bump is \bnone{} although an older minor or major change
established the train; read literally, the summary loses the aggregate the
specific rule requires. This report follows \spec{11.5}; the summary should
delegate to it or carry a separately named train target.
\end{obligation}

\paragraph{What changed since 2.0.0.} The obligation as first recorded had a
second half: \spec{11.5} said \diag{E185} fires when the target is not
greater than the ``baseline core'', which would have refused every ordinary
graduation. The current text resolves it: ``Graduation never lowers a
version: if \code{target} is lower than the core of the baseline, \code{E185}
is raised. A \code{target} equal to that core is the ordinary graduation,
because \code{1.3.0} ranks above \code{1.3.0-rc.1}.'' The same section now
also carries the channel-entry patch on the way out of a train (\diag{W204}
on graduation, vector 39b), which removes the dead end in which a train
entered by the patch could not graduate. Both are implemented. The first
half of the obligation stands: \spec{13.6}'s own second paragraph says that
``\S11.4 then uses the full \code{W(P)} to compute that admitted
prerelease's aggregate target'', and the window table of \spec{13.3} assigns
the aggregate to $W(P)$, but \spec{13.9} is the section an implementer
transcribes and it names neither.

\paragraph{Consequence.} \Cref{thm:g3} proves that fixed version inputs give
the same result on retry; it does not decide between two formulas.
\Cref{def:plan} follows \spec{11.5}. A second implementation transcribing
\spec{13.9} alone graduates a channel-only train to its stable baseline,
which \spec{13.9}'s own last guard then refuses as \diag{E195}.

\paragraph{Pin.} The last row of \spec{11.4}'s table (``\code{1.3.0-rc.1};
nothing new, \code{E} is still the window's \code{minor}; \code{stable};
\code{1.3.0}; graduation'') and edge case 61 of \spec{15.5}: ``Graduation
with no pending bumps: publishes the accumulated \code{target}.'' The engine
follows the specific rule: \code{computeVersion} in \code{plan.go} applies
the bump held over the whole train to the stable baseline, and
\code{TestTrainGraduationCountsPublishedWork} asserts that ``the bump of the
already-published breaking change must still be in force even though it
alone would not have re-released the train'';
\code{TestGraduationUpdatesSpanTheTrain} and
\code{TestTrainEnteredByTheEntryPatchGraduates} pin the record and the
patch.

\subsubsection{One edge-kind set orders, blocks and refuses cycles}

\begin{obligation}[no \code{orderKinds} or \code{blockingKinds} in the engine]
\label{ob:kinds}
\spec{13.1} requires the graph to be acyclic ``over the edge kinds of
\code{propagation.kinds} and \code{publish.orderKinds}'', and adds that a
cycle ``existing only through [\code{devDependencies}] is not a cycle for
this purpose and does not trigger \code{E200}, which is the common and
legitimate case''. \spec{19.2} orders over \code{publish.orderKinds},
``which includes \code{optionalDependencies}''; \spec{19.3} blocks over
\code{publish.blockingKinds}, ``which excludes \code{optionalDependencies}''.
The engine has neither key. Every declared edge, whatever its kind, enters
the one graph that the cycle check, the publish order and the failure
closure read; only propagation honours a kind set.
\end{obligation}

\paragraph{Where.} \code{config.go} builds one \code{model.Dependency} per
declared edge, carrying its kind and dropping none; \code{loadWorkspace} in
\code{plan.go} adds each to the graph and reports \diag{E200} on any cycle
of it; the publish order is that graph's topological order; and the blocking
closure of \code{release/executor.go} reads a per-relation boolean,
\code{isBlocking}, not the edge's kind. The kind set of
\code{parser.propagation.kinds} is consulted by the propagation walk alone,
which is what the integration test plan states as the purpose of
\code{TestPropagationKindsSelectTheEdgesTraversed}. Neither the design
history, the changelog, the release notes nor the documentation records the
difference; the documentation of \code{propagation.kinds} describes the walk
only.

\paragraph{Scenario.} Two packages that list each other under
\code{devDependencies}, a test fixture depending back on the package it
exercises. The specification accepts the workspace (vector 100a, edge case
99a); the engine refuses it with \diag{E200} before any plan. Separately,
\code{docs-site} depending on \code{ui} under \code{optionalDependencies}
with \code{ui}'s publish failing: the specification publishes
\code{docs-site} (vector 79, edge case 100); the engine blocks it with
\diag{W194} unless its relation declares \code{isBlocking: false}.

\paragraph{Consequence.} The departure is on the safe side. The engine's
graph is a superset of the one \spec{19.2} orders, so any topological order
of it is an order the specification accepts and \cref{prop:order} is
unaffected; the blocked set of \cref{def:blocking} is computed over a
superset of \code{blockingKinds}, so it is a superset of the specified one
and the $F \cup B$ of \cref{thm:resumption} is larger, never smaller.
Nothing the specification forbids is published. What is lost is admission: a
workspace with a development-only cycle cannot release at all, and an
optional dependent waits for a provider it does not need.

\paragraph{Repair.} Either implement the two keys with the specified
defaults, or record the departure where the design history records the
others and restrict \diag{E200} to the runtime kinds. The second is a
one-line change to the graph the cycle check reads.

\paragraph{Pin.} Vectors 79 and 100a and edge cases 99a and 100 are the
conformance cases; none has a test, and both vectors fail against the engine
as it stands. \code{TestFatalDependencyCycle} of the integration suite and
\code{TestCycleIsRejected} of \code{internal/plan} pin \diag{E200} over
runtime edges; \code{TestExecutionDevDependencyOutputReachesConsumer} shows a
\code{devDependencies} edge ordering builds under \spec{28}, which is what
vector 22 of \spec{28.8} asks.

\subsubsection{The fleet recorder rejects a retry that annotates the same commit}

\begin{obligation}[the create-only rule's exception is implemented on one of two paths]
\label{ob:records}
\spec{19.1} states: ``A release tag is created, never replaced. Writing it
to the authoritative store MUST fail when the store already holds that tag
name, unless the stored tag names the same commit for the same package and
version, which is the retry of an uncertain write (\S19.4).'' The
single-repository recorder implements both halves; the fleet recorder of
\spec{27} implements the first and not the exception.
\end{obligation}

\paragraph{Where.} \code{gitx/records.go} leases every record on the name
being absent and, when the store refuses, reads the stored tag back peeled:
the same commit is ``this run's own write arriving twice and is a success'',
another commit is a published record left where it is (\diag{E221}). Its
inventory compares peeled lines because ``two runs annotating the same
commit write different tag objects, and a comparison on those would call one
release two''. The fleet recorder's \code{PushRelease} in
\code{gitx/polyrepo\_write.go} pushes each source tag with a plain unforced
\code{git push}, which Git rejects whenever the remote holds a different tag
object under that name, the same commit included; the rejection surfaces as
\diag{E335} after the publication.

\paragraph{Scenario.} One run's tag push is applied by the remote and its
acknowledgement is lost, or two checkouts of one source repository record
the same version at the same head in turn. The next push of that name
carries a new annotation of the same commit. In a single repository the push
is reported as \diag{W223}, ``tag already exists on the remote at this
release's commit, skipped''; in a fleet it is refused, the run reports
\diag{E335}, and the dependent legs are blocked (\spec{27.7}).

\paragraph{Consequence.} Nothing incorrect is recorded: the refused write is
the identity the store already holds, and the next run reads it as released.
The cost is a failed run and blocked consumers for a retry the specification
defines as successful. \Cref{cor:exactly-once} is unaffected, since the
ledger holds one record either way.

\paragraph{Status.} Open for the fleet path; resolved for the
single-repository path by the design history entry of 21 September 2026,
``Record authority under the lock''. Neither the design history, the
changelog nor the test plan records the fleet path's difference.
\code{TestReleaseRecordsCreateOnlyPush} pins the single-repository
behaviour (``the same commit is the retry of a lost write and succeeds'');
\code{TestReleaseRecordsInAComposedWorkspace} pins the fleet's comparison
under the lock, not its push. Cases: \spec{28.8} vectors 27 and 28; edge case
76a.

\subsubsection{The failure category travels only from a distributed lock loss}

\begin{obligation}[the class of \spec{28.9} is attached at the new sites only]
\label{ob:category}
\spec{28.9} requires implementations to ``expose stable machine-readable
categories, named by the identifier each row states''. The design history
entry of 22 September 2026 states the engine's position: ``the existing
\code{E220} to \code{E222}, \code{E335} and \code{E336} carry
\code{native-recording-or-lock} as their class; the field is attached where
a distributed run loses its lock, and not yet at the older recording and
lock sites, which is a departure to be closed by attaching it there.''
\end{obligation}

\paragraph{Consequence.} A reader that switches on the class, which is what
\spec{28.9} designs the class for, sees a lost lock in a distributed run and
not the same lost lock in a local one, nor a refused record. No release
outcome depends on it. \code{execution/codes.go} names the mapping in one
place, so the repair is to emit the field at the five older sites.

\paragraph{Pin.} Vector 14 of \spec{28.8} and
\code{TestExecutionLockLossStopsNewEffects} of the integration suite cover
the site that carries the field; row EXE-12 of the requirement matrix
asserts \diag{E226} with its \code{execution-authority} class. No test
asserts the class on a local \diag{E221} or \diag{E336}.

\subsubsection{Profile departures and inapplicable vectors}

\begin{obligation}[what the engine records as not implemented or done otherwise]
\label{ob:profile}
The design history entry of 22 September 2026 lists, beside
\cref{ob:category}, the remaining departures of the engine from \spec{28}
and the vectors that do not apply to it. They are recorded here so that the
table of \cref{sec:conformance} reads against them.
\end{obligation}

\begin{itemize}
  \item ``Not implemented, so the conditional vectors do not apply: reuse of
    verified outputs across runs (vector 20) and the rollback profile of
    section 26 (vector 26).'' \spec{28.8} states the rule the entry relies
    on: a vector whose condition does not hold for an implementation ``is
    neither satisfied nor violated''.
    \code{TestExecutionOldRunReceiptGrantsNothing} and
    \code{TestExecutionTaskReceiptDischargesNothing} pin the halves of those
    vectors that hold regardless.
  \item ``Departures from the profile as written: the platform check of
    section 28.2 runs before dispatch for the packages a run releases and,
    for a prepared provider, at its placement, where an unsatisfiable
    platform fails the provider's consumers.'' \spec{28.2} asks that
    unsupported target platforms ``fail preflight''.
    \code{TestExecutionPreparedProviderObeysItsOwnPlatforms} and
    \code{TestExecutionOrchestratorOnlyBuildIsNotHeldToTheWorkersPlatforms}
    pin the engine's timing.
  \item \spec{28.4} ``now also permits a durable mailbox ref or a
    history-free snapshot, and neither is implemented''. These are
    permissions, not requirements; the entry records the cost of the mode in
    use, a first transport of a whole history.
  \item The specification's own header keeps the VCS adapters of \spec{25}
    and the explicit rollback of \spec{26} ``pending implementation''. The
    forward-publish projection this report models is unaffected;
    \spec{13.7c} says its guarantees are claims about that projection.
\end{itemize}

\subsubsection{Textual defects noticed in passing}

Four defects of the text, none of which changes a rule, are recorded for the
specification's owner because they affect citation. Appendix B numbers two
vectors ``82c'': the consumer that gets ahead of a held provider (B.7) and
the hoisting of \code{resolvableBy} over a mixed source set (B.7a); this
report cites them as ``82c (B.7)'' and ``82c (B.7a)''. In \spec{19.3} the
bullets ``continue publishing packages that are not blocked'' and ``not
release a package at the baseline commit of a consumer it still owes'' are
run into the preceding bullet's line rather than set as items. The design
history says the release notes list the departure of \cref{ob:owed}, and at
this commit they do not. And the window table of \spec{13.3} cites
\spec{13.6} as a consumer of the aggregate train window, while the formula
of \spec{13.6} reads the undelivered one and only its prose defers the
aggregate to \spec{11.4}, the seam \cref{ob:graduation} sits on.

\subsection{Historical CCME 1.0.0 obligations}
\label{sec:obligations-historical}

Formalising CCME 1.0.0 at \texttt{713f1a0b} exposed six failure classes. Each obligation below
states the claim, a minimal counterexample, the consequences, a repair, and a
conformance vector. Three concern \spec{13.7c}, one \spec{13.7b}, one
\spec{9.3}, and one the optimisation advice of \spec{13.11}.

None is a defect in the reference implementation. In two cases
(\cref{ob:pre,ob:duality}), the implementation already uses a mechanism absent
from the specification. A second implementation following only the written
contract would therefore be incorrect. Source references below are
relative to \code{services/dispat/internal/plan/} and cite the line numbers of
the released \code{services/dispat/v1.7.2} tag. Those sources have since
moved: at \code{252216ba}, \code{containedInBaseline} is at \code{plan.go}
line 2507, \code{Plan.StaleSources} at line 1280, and the per-unit traversal
of \code{propagate.go} is served from a walk cache (line 76) whose
composition is checked against the literal walk (\cref{sec:conformance}).

Released CCME 2.0.0 resolves them as follows: \cref{ob:pre} by $W_{\rm fresh}$;
\cref{ob:g5,ob:origin} by the full retry invariant and fixed-state wording;
\cref{ob:duality} by applying channel eligibility in the upward audit;
\cref{ob:g8} by stating axis independence in the bump-to-channel direction
only; and \cref{ob:bucket} by retaining each unit's own source exclusion when
bucketed. The counterexamples remain regression evidence for those normative
conditions. The regression vectors named below as ``the released 2.0.0
form'', 133 to 138, are in the current text the edge cases 133 to 138 of
\spec{15.7}, one row each. Except where a paragraph explicitly names CCME 2.0.0 or a current
theorem, every numbered \spec{} reference and quoted guarantee in this subsection
refers to the CCME 1.0.0 snapshot at \code{713f1a0b}; present tense describes
the historical rule being executed, not the released 2.0.0 text.

\subsubsection{The prerelease window never discharges}

\begin{obligation}[the 1.0.0 G4/G6 claims fail on a prerelease train]
\label{ob:pre}
Under the 1.0.0 single-window rule, ledger discharge and convergence fail after a successful prerelease publish. Current \cref{thm:g4,thm:g6} use $W_{\rm fresh}$ and do not make that claim.
\end{obligation}

\paragraph{Counterexample (bump axis).}
This is \SPEC{} vector 95b, run twice. Baselines
\code{core@1.5.0-beta.0} and \code{cli@2.0.1-beta.0}, both with an older stable
baseline; \code{cli} depends on \code{core}; one commit
\code{fix(core)\^{}: x}.

Run 1 plans \code{core@1.5.0-beta.1} and \code{cli@2.0.1-beta.1} and publishes
both. Run 2, at the same \Head{}: by \cref{lem:window} neither window changed,
because \cref{def:window} measures from the last \emph{stable} tag and no
stable tag was written. Under the 1.0.0 direct and propagated rules, \code{core} and \code{cli} each receive \bpatch{} again, both packages are planned again, and
\spec{11.4} increments the counters to \code{beta.2}. Run 3 gives
\code{beta.3}. The plan never empties, and \diag{E199}, the convergence
check of \spec{19.6}, which admits no exclusion but held packages, fires
after the first fully successful run of any prerelease train.

\paragraph{Counterexample (channel axis).}
\Cref{thm:g7} says the channel a package arrived on is not proposed for it
again. It does not say nothing is proposed. Let $p$ have stable baseline
\code{1.2.0} and a window containing an older \code{release(p)\%rc} and a
newer \code{release(p)\%beta}. Run 1: \spec{11.6} precedence gives the newer
directive, $p$ is released \code{1.2.1-beta.0} (\diag{W186} noting the
conflict), and the operator's stated intent is honoured. Run 2: the window is
unchanged, so both directives are still candidates; \code{\%beta} now proposes
nothing (\diag{W199}, \cref{thm:g7}) and the \emph{older} \code{\%rc},
which \spec{11.6} had ruled had lost, wins by default and $p$ is released
\code{1.2.1-rc.0}. Run 3: \code{\%rc} is now inert and \code{\%beta} proposes
\code{beta} again, computing \code{1.2.1-beta.0}, which does not exceed the
baseline \code{1.2.1-rc.0}, so the run stops with \diag{E195}. A repository
that did nothing wrong is wedged, having published a version onto a channel its
newest directive did not ask for.

\paragraph{What broke in 1.0.0.} Its G4 (no double release), G6 (convergence), and G7 channel-settling claim (the channel axis settles) all fail, and with
them \spec{17.2}'s idempotency clause and
\spec{19.6}'s convergence check, for every package on a prerelease train,
which is to say for the entire feature the prerelease flow exists to support.

\paragraph{The repair.} The window answers one question, what has $p$ not
yet released, and \spec{13.3} correctly measures it from the stable tag
because \spec{11.4} recomputes a train's target over the whole train. The
missing distinction is that the same set is then used to answer a second
question: is anything still \emph{owed}? Separate them. Define the
\emph{fresh window}
\[
  \Winf{\st}{p} \;=\; \Win{\st}{p} \setminus \Rch{\basec{\st}{p}},
\]
which for a package whose baseline is stable equals $\Win{\st}{p}$, and for one
on a train excludes the commits its own prereleases already published. Then:
\begin{enumerate}[label=(R\arabic*),leftmargin=3.2em]
  \item \cref{def:sources} discharges on $\Winf{\st}{p}$, so a package whose
    prerelease shipped the work counts as a source (which \spec{13.4a} already
    intends: the artefact is public);
  \item plan membership (\cref{def:plan}) requires a contribution from a commit
    in $\Winf{\st}{p}$, while the version computation (\spec{11.4}) keeps
    reading $\Win{\st}{p}$, so the train's target still accumulates over the
    whole train;
  \item candidate collection in \cref{def:phase1,def:channels} is restricted to
    $\Winf{\st}{p}$, which is what stops the superseded directive of the second
    counterexample from voting again.
\end{enumerate}
With (R1)--(R3), \cref{lem:window} strengthens to $\Winf{\st'}{p} = \emptyset$
for \emph{every} released $p$, prerelease or not, because $\basec{\st'}{p} =
\Head$. Fresh discharge repairs the prerelease counterexample. Current \cref{thm:g4,thm:g5,thm:g6} separately state the full retry-invariant conditions needed for persistence, non-widening, and bounded convergence; (R1)--(R3) alone do not imply them.

\paragraph{Evidence.} The reference implementation has exactly this predicate,
under a different name and undocumented in \SPEC{}:
\code{containedInBaseline} (\code{plan.go},
lines 1836--1855) reports whether a commit has already been published by the
package's baseline tag, and its callers implement (R1) in
\code{sourcePackages}, (R2) by gating \code{NewWork} on it, and (R3) in both
\code{resolveChannels} and \code{propagateChannels}. Its comment states the
consequence the specification does not: without it the engine ``would re-warn
\dots on every run until graduation''.

\paragraph{Regression vectors.} Vector 133 is the released 2.0.0 form of the bump-axis case. Historically, (i) vector 95b, run twice: the second run's plan
must be empty. (ii) The two-directive window above, run three times: run 2 must
be empty, and \diag{E195} must not arise.

\subsubsection{Window-scoped suppression lapses, so the admitted set can grow}

\begin{obligation}[the 1.0.0 G5 claim lacks suppression stability]
\label{ob:g5}
The 1.0.0 proof assumed ``same sources, same depth, same graph at \Head{}'', although window-scoped suppression can change the source set. CCME 2.0.0 makes the needed stability part of the full retry invariant used by current \cref{thm:g5}.
\end{obligation}

The root cause is a design decision taken three times, in three sections, each
time for a good local reason: suppression is scoped to the window. A
\code{cancel} stops counting once the package releases (\spec{13.4a}); a hold
does the same (\spec{13.4a}); and a correction ``reaches only undischarged
work'', so a \code{Deletes} is a no-op for a package that has released its
target (\spec{7.4.2}, \cref{rem:corrections}). Each rule is right on its own
terms. Together they mean that the set of units a package acts as a propagation
source for is \emph{monotone in the state}, which is precisely what \G{5} says
it is not.

\paragraph{Counterexample (cancellation).} \code{cli} depends on \code{core}.
Three commits: $A$ carrying \code{feat(core)\^{}: streaming reader}, then $B$
carrying \code{cancel(core)}, then $C$ carrying \code{fix(core): unrelated}, in
that ancestry.

Run 1: $\canc{A}{\code{core}}$ holds and $A \in \Win{\st_1}{\code{core}}$, so
\code{core} is not in $\srcp{\st_1}{u_A}$ and $u_A$ propagates to nothing;
\code{core} releases \code{1.4.3} from $C$ alone (\SPEC{} vector 41).
$\plan{\st_1} = \{\code{core}\}$: \code{cli} is not planned.

Run 2: \code{core} is tagged at \Head{}, so $\Win{\st_2}{\code{core}} =
\emptyset$, so $\disch{\st_2}{\code{core}}{A}$ holds, so \code{core} re-enters
$\srcp{\st_2}{u_A}$ by the discharge disjunct of \cref{def:sources}. Now $u_A$
propagates: $A \in \Win{\st_2}{\code{cli}}$, \code{cli} is not cancelled, and
the bump is admitted. $\plan{\st_2} = \{\code{cli}\} \not\subseteq
\plan{\st_1}$.

\paragraph{Counterexample (corrections).} The same shape without any
\code{cancel}. Let $B$ instead carry \code{chore(core): drop it} with
\code{Deletes: <A>}. In run 1 the correction discards $A$'s record for
\code{core}, so $u_A$ has no source and reaches nothing; \code{core} still
releases, on $C$. In run 2, $\Win{\st_2}{\code{core}}$ is empty, so by
\spec{7.4.2} the correction is a discharged no-op for \code{core}
(\diag{W209}), $A$'s record stands again, and $u_A$ propagates to \code{cli}
exactly as before. This route is worth stating separately because \diag{W209}
reads as reassurance (the correction did not need to do anything) when it
is in fact the moment the suppressed propagation returns.

\paragraph{What broke in 1.0.0.} Its unrestricted G5 claim and the ``nothing is added'' step of its G6 proof failed. \Cref{thm:g6}'s conclusion survives here (\code{cli}
releases in run 2 and run 3 is empty) but the variant argument does not, and
the security reading of \G{5} does: whether \code{cli} is released at all
depends on whether \code{core}'s publish succeeded. Inducing a publish failure
changes the set of packages a commit eventually releases, which is the
inference \spec{18.1} relies on \G{5} to forbid.

\paragraph{The repair, and why the obvious one is wrong.} The tempting fix is
to make cancellation sticky, never letting a cancelled package become a source,
by defining $\disch{\st}{p}{c} \land \lnot\canc{c}{p}$. It is wrong: it
breaks \SPEC{} vector 72, where \code{cancel(core)} lands \emph{after}
\code{core} published the work and the consumers' catch-up must still proceed.
Ancestry cannot separate the two cases, since in both the cancel is in the
history at \Head{} and in both \code{core} carries a tag at \Head{}. The
difference is only whether the released version counted the unit, and that is
not recoverable from the state.

The recommendation is therefore to restate the guarantee rather than change the
algorithm: \G{5} holds for every unit none of whose sources is suppressed at
\Head{}, which is hypothesis (H2), and the widening in the remaining case is
intended behaviour, since neither a \code{cancel} nor a \code{Deletes} touches
code (\spec{10.1}, \spec{7.4.2}) and the work of $A$ is in fact inside the
published \code{core@1.4.3}. What the specification should not do is claim a
subset property that the discharge rule of \spec{13.4a} deliberately violates,
and \spec{18.1} should not rest a security argument on the unrestricted form.

\paragraph{Evidence.} The cancellation counterexample was run against the
reference implementation, which follows \spec{13.4a} literally and reproduces
it: with the history above and \code{core} tagged at \Head{} in the second
state, the first plan releases \code{core} alone and the second releases the
consumer as well, reporting it as a \emph{catch-up from core}. The engine is
doing what the specification says; it is \G{5} that does not describe it.

\paragraph{Regression vector.} Vector 134 is the released 2.0.0 form. In the historical pair, each history above is run twice: run 2
plans \code{cli}, the plan is documented as a superset of run 1's, and the
\diag{W193} naming \code{core} as the origin is present so that the widening is
visible to whoever reads the plan rather than silent.

\subsubsection{The staleness dual omits the resolvability test}

\begin{obligation}[\spec{13.7b}'s two formulations disagree]
\label{ob:duality}
\spec{13.7b} states that \code{staleSources(D)} ``is non-empty exactly when
\S9.2 assigns \code{D} a non-\code{none} propagated bump'', and that
disagreement is an implementation bug. The procedure tests scope, depth,
window, and cancellation, but has no counterpart to \cref{a:resolvable}.
\end{obligation}

\paragraph{Counterexample.} \SPEC{} vector 95 (also edge case 60m):
\code{feat(core)\^{}\%beta: x} where \code{core} moves onto \code{beta} and its
consumers stay on \code{stable}. \spec{9.3a} suppresses every consumer's bump
with \diag{W208}, so $\propb{\st}{\code{ui}} = \bnone$, and \SPEC{} is
explicit that ``\code{feat(core)\^{}\%beta} releases \code{core} alone''. But
\code{staleSources(ui)} passes every guard it has and returns the row
$(\code{core}, u, 1, \bpatch)$. The audit reports \code{ui} as behind a
dependency for which no release is owed, and by \spec{13.7b}'s own words the
implementation is now buggy by definition.

\paragraph{The repair.} Add the hoisted predicate of \cref{prop:hoist} to the
procedure, immediately after the \code{cancelledFor} test:
\begin{speccode}
        srcChan = { channel[P] : P in sources }        # §9.3a, as in §9.2
        if not ('stable' in srcChan or channel[D] in srcChan):  continue
\end{speccode}
With it, \cref{thm:duality} holds unconditionally and its hypothesis can be
dropped.

\paragraph{Evidence.} The reference implementation cannot exhibit the
disagreement, because \code{Plan.StaleSources}
(\code{plan.go}, line 1217) returns the rows
recorded by the downward traversal rather than re-deriving them from an
independent upward walk. The duality is therefore true there by construction
and untested; an implementation that follows \spec{13.7b} literally, as a
second implementation would, disagrees on every prerelease workspace.

\paragraph{Regression vector.} Vector 135 is the released 2.0.0 form; historically, vector 95 was extended: assert that
\code{staleSources(ui)} is empty, and that the randomised
down-versus-up cross-check \spec{13.7b} recommends is run on a workspace
containing at least one prerelease train.

\subsubsection{Axis independence is one-directional}

\begin{obligation}[the ``vice versa'' in \G{8}]
\label{ob:g8}
\spec{13.7c} states \G{8} as ``suppressing a bump under \spec{9.3a} never
suppresses a channel, and vice versa'', and proves it by observing that the two
passes share only $\chanres{\st}{\cdot}$, which phase 3 reads and never writes. That
argument proves one direction only, and the other is false.
\end{obligation}

\paragraph{Counterexample.} Let \code{core} and its consumer \code{app} both
sit on \stable, and let one commit carry \code{feat(core)\^{}\%beta++1}: the
caret propagates a bump to \code{app} and the \code{++1} propagates the
\code{beta} channel to it, which is exactly the pairing \spec{9.3a} prescribes
for taking consumers along. Suppose \code{app}'s own window also contains
\code{release(app)\%rc}. By \cref{def:channels} a direct directive beats every
propagated one, so $\chanres{\st}{\code{app}} = \code{rc}$ and the propagated
\code{beta} is discarded. Phase 3 then computes
$\srcch{\st}{u} = \{\code{beta}\}$, finds
$\chanres{\st}{\code{app}} = \code{rc} \notin \srcch{\st}{u}$ and no source on
\stable, and suppresses the bump with \diag{W208}. Had the propagated channel
been applied, \code{app} would have been on \code{beta} and the bump admitted.
The suppression of a channel proposal caused the suppression of a bump.

\paragraph{The repair.} State \G{8} as \cref{thm:g8} does: the channel
assignment is computed by phases 1--2 and reads no phase-3 output, so no bump
suppression can affect it. The dependency in the other direction is real, is
the entire content of \spec{9.3a}, and should be documented as such rather than
denied. The operational consequences \spec{13.7c} draws from \G{8} are
unaffected: they only ever use the direction that holds.

\paragraph{Regression vector.} Vector 136 is the released 2.0.0 form. The historical workspace above asserts that \code{app} is
released on \code{rc} with no propagated bump and a \diag{W208}, and that the
\diag{W208} is a consequence of the channel precedence rule rather than of
anything the caret did.

\subsubsection{Inherited channels read a baseline that release moves}

\begin{obligation}[\spec{9.3}'s replay-independence claim]
\label{ob:origin}
\spec{9.3} states of the \code{inherit} resolution that ``the result is
deterministic, does not depend on which earlier run published the origin''.
The second half is false: resolution rule 2 reads
$\chanof{\base{\st}{q}}$, and the origin's own release moves it.
\end{obligation}

\paragraph{Counterexample.} \code{app} depends on \code{core}; \code{core}
sits at baseline \code{1.5.0-beta.0} over stable baseline \code{1.4.2};
\code{app} at \code{2.0.0}. Two commits: $C_1$ carrying
\code{feat(core)\^{}++1: x} (no explicit channel value, so the propagated
channel is \code{inherit}) and $C_2$ carrying
\code{release(core)\%stable: graduate}.

Run 1: \code{core}'s own channel is \stable{} (the direct directive), but
the inherited value for \code{app} is
$\chanof{\base{\st_1}{\code{core}}} = \code{beta}$, so \code{app} is
planned at \code{2.0.1-beta.0} while its provider graduates to
\code{1.5.0}. Suppose \code{core} publishes and \code{app} fails. Run 2:
$\chanof{\base{\st_2}{\code{core}}} = \stable$, the proposal is inert
(\diag{W199}), and \code{app} is planned at \code{2.0.1}. Both runs plan
\code{app}, at different versions, with every completed release stable:
the 1.0.0 cross-run version-stability claim fails. CCME 2.0.0 instead fixes inherited baseline-channel state as one component of its full retry invariant, as required by current \cref{thm:g3}. Run against the reference implementation, the two
plans are exactly as computed here (\code{app}: \code{beta}
$\rightarrow$ \code{stable}, \code{2.0.1-beta.0} $\rightarrow$
\code{2.0.1}).

\paragraph{What breaks, and what does not.} Determinism
(\cref{thm:determinism}) is unaffected: each plan is still a pure function
of its state. What fails is the \emph{cross-run} half of the specification's
sentence, and with it channel stability for unreleased dependents of a
graduating origin. Note also that run 1's own output is operationally odd:
a consumer is dragged onto a train in the same run that ends it. Both
effects have one cause, phase 1 reading baselines so that it cannot depend
on phase 2's output (\spec{9.2}); the price of that layering is that
\code{inherit} answers with the origin's pre-graduation channel.

\paragraph{The repair.} Scope the specification's sentence to a fixed state
(determinism per run, which is what \cref{thm:determinism} proves), and
state the drift explicitly. Resolving \code{inherit} against the origin's
\emph{resolved} channel instead would remove the drift but requires a
fourth phase, which \spec{9.2}'s no-circularity argument was written to
avoid; the fresh window of \cref{ob:pre} does not remove it either, since
the commit remains pending for the failed dependent. The honest fix is
documentation plus the vector below.

\paragraph{Regression vector.} Vector 137 is the released 2.0.0 form: the two commits above, two runs: run 2 must
plan \code{app} at \code{2.0.1} on \stable{}, and the plan should carry the
\diag{W199} that records why the run-1 proposal lapsed.

\subsubsection{Bucketed traversal over-admits a unit's own sources}

\begin{obligation}[the 1.0.0 bucketing remedy has its own edge]
\label{ob:bucket}
CCME 1.0.0 \spec{13.11} warned that subtracting the merged source set from a bucketed
traversal silently withholds bumps, and prescribes admitting ``any node reached
across at least one edge within the depth bound''. By \cref{prop:bucket} that
prescription over-admits the set $X$: a package that is itself a source of the
unit that reaches it.
\end{obligation}

\paragraph{Counterexample.} One unit, \code{fix(core,cli)\^{}major: x}, where
\code{cli} depends on \code{core}. \spec{9.2} seeds \code{seen} with both
sources, so \code{cli} is not a propagation target and takes the unit's own
\bpatch. Under the merged rule, \code{cli} is reached from \code{core} across
one edge inside the depth bound and takes the propagated \bmajor{} as well:
a two-level difference in the published version, with no diagnostic.

\paragraph{The repair.} Bucket the traversal, then re-decide the few targets
that are themselves sources. Since the bucket key fixes the propagated bump,
the only question for such a $d$ is whether \emph{some} unit in the bucket
admits it, that is whether $\exists i: d \notin S_i \land \dist{S_i}{d} \leq
\ell$. The candidates are confined to $S$, which is small, and the per-unit
source sets are already materialised, so the check is cheap. Equivalently, and
more simply, exclude from bucketing any unit whose own source set contains a
package reachable from that set within the depth bound.

\paragraph{Evidence.} The reference implementation does not bucket: it
traverses once per unit (\code{propagate.go},
line 401), so it is unaffected. The obligation is on the advice, which a
performance-minded second implementation would follow.

\paragraph{Regression vector.} Vector 138 is the released 2.0.0 form. Historically, \SPEC{} Appendix B.7a, whose vectors test
exactly this class of almost-equivalence, gains one: \code{fix(core,cli)\^{}major}
must publish \code{cli} as a patch, bucketed or not.
